\section{Performance space and the preferred direction solver}
\label{sec:performance}

The preceding coordinates modify current space. A more radical approach maps
each direction directly to what it accomplishes. In whitened coordinates let
\(\vect x=\gls{sym:y}/\norm{\gls{sym:y}}\), so \(\norm{\vect x}=1\), and set
\begin{equation}
 \bar{\matr H}=\matr R^{-1/2}\matr H\matr R^{-1/2},\qquad
 \bar{\matr Q}=\matr R^{-1/2}\matr Q_U\matr R^{-1/2}.
\end{equation}
Define
\begin{equation}
 \gls{sym:tau}(\vect x)=\vect x^{\mathsf T}\bar{\matr H}\vect x
 =\frac{M}{P_{\mathrm{Cu}}},\qquad
 \gls{sym:nu}(\vect x)=\vect x^{\mathsf T}\bar{\matr Q}\vect x
 =\frac{U^2}{P_{\mathrm{Cu}}}.
 \label{eq:performance-coordinates}
\end{equation}
Both are invariant under nonzero current scaling. The horizontal coordinate
is squared voltage consumed per watt of copper loss; the vertical coordinate
is torque produced per watt of copper loss.

For two real symmetric forms and three or more real variables, the image of
the unit sphere
\begin{equation}
 \mathcal W=\{(\nu(\vect x),\tau(\vect x)):\norm{\vect x}=1\}
 \label{eq:joint-range}
\end{equation}
is compact and convex. Brickman's unit-sphere result is the exact statement
needed here \cite{brickman1961}; Dines's related theorem says that the
unnormalized joint range of two homogeneous quadratic forms is a convex cone
\cite{dines1941}. The dimension qualification matters. In two real variables,
the sphere is a circle and its joint range can be only the boundary of an
ellipse rather than its filled interior. The extra EESM current coordinate is
what makes the normalized real performance image convex.

\begin{figure}[tb]
\centering
\begin{tikzpicture}
\begin{axis}[axis main,width=.8\linewidth,height=6.2cm,
 xlabel={$\nu=U^2/P_{\mathrm{Cu}}$},ylabel={$\tau=M/P_{\mathrm{Cu}}$},
 xmin=0,xmax=5.2,ymin=-2.1,ymax=2.8]
 \addplot[draw=ForestGreen!70!black,very thick,fill=ForestGreen!15,
 domain=0:360,samples=161]
 ({2.4+1.75*cos(x)},{.25+1.35*cos(x)*.25+1.35*sin(x)*.97});
 \addplot[support line] coordinates {(.5,2.55) (4.8,1.25)};
 \addplot[voltage curve] coordinates {(0,0) (4.9,2.45)};
 \addplot[only marks,mark=*,BrickRed] coordinates {(1.55,2.23)};
 \node[annotation] at (axis cs:3.9,2.0){voltage-feasibility line};
 \node[annotation] at (axis cs:1.55,-1.45){compact convex direction image};
\end{axis}
\end{tikzpicture}
\caption{Every point represents a current direction. A rotating support line
finds the upper boundary; the requested-torque voltage inequality is a line
through the origin. Their contact selects the optimal direction.}
\label{fig:performance-space}
\end{figure}


The boundary can be found without sampling the sphere. Rotate a support line
with normal \((-\mu,1)\). Its highest contact maximizes
\begin{equation}
 \tau-\mu\nu
 =\vect x^{\mathsf T}(\bar{\matr H}-\mu\bar{\matr Q})\vect x.
\end{equation}
On the unit sphere, the maximizing direction is the dominant eigenvector of
the symmetric pencil
\begin{equation}
 (\matr H-\mu\matr Q_U)\vect v
 =\lambda\matr R\vect v.
 \label{eq:support-pencil}
\end{equation}
This is not a disguised request to memorize a multiplier. Geometrically, the
controller turns a support line in a two-dimensional convex performance set;
the eigenvector is the current direction at the contact point.

With radial loss scale \(P\), every performance point obeys
\begin{equation}
 M=P\tau,\qquad U^2=P\nu.
\end{equation}
For requested positive torque, the admissible directions satisfy
\begin{equation}
 \nu\leq\frac{U_{\max}^2}{M_{\mathrm{ref}}}\tau,\qquad \tau>0,
 \label{eq:performance-wedge}
\end{equation}
and the best point is the one with largest \(\tau\) in this half-plane.
Therefore
\begin{equation}
 \boxed{P_{\mathrm{req}}=\frac{M_{\mathrm{ref}}}{\tau_\star}.}
 \label{eq:performance-preq}
\end{equation}
If the unconstrained topmost point already satisfies the sloped voltage line,
one dominant eigenvector of \(\bar{\matr H}\) suffices. Otherwise the optimum
lies where the performance boundary meets that line. A scalar search in
support slope \(\mu\) locates the contact.

For maximum torque, one direction supplies
\begin{equation}
 M_{\mathrm{dir}}=\tau\min\!\left(P_{\mathrm{Cu},\max},
 \frac{U_{\max}^2}{\nu}\right),
 \label{eq:performance-mmax}
\end{equation}
with \(U_{\max}^2/\nu=+\infty\) when \(\nu=0\). The maximum along the upper
boundary of \(\mathcal W\) is found by the same support-line parameterization.
No 3D gridding is required.

\subsection{SVD, GSVD, and what each axis optimizes}

The SVD of \(\matr A\matr R^{-1/2}\) diagonalizes voltage gain per loss and
reveals the voltage-null direction robustly. The eigenvectors of
\(\bar{\matr H}\) diagonalize torque per loss. A generalized SVD can expose
relative voltage/loss gains, but it does not generally diagonalize torque as
well. Performance space avoids choosing one permanent axis system: it uses a
small pencil whose axes change continuously with the support slope. This
runtime flexibility is the reason it is preferred over a single simultaneous
diagonalization.
